% ---------------------------------------------------------------------------
% Section II Related Work, shortened 2026-09-13 for the page budget from the
% mentor's version of 2026-09 (1,002 words, 1.2 pages) to about 620 words.
% Every citation key of the mentor's version is kept. Changes: subsection B
% (Static and Conditional Interventions) folded into one paragraph at the end
% of subsection A because Section III-A makes the same distinction; the
% Takeaway box removed because the closing paragraph of subsection C says the
% same; the remaining paragraphs tightened without removing any claim or
% number. Mentor-owned section, offered as a replacement.
% ---------------------------------------------------------------------------

\section{Related Work}

\subsection{Training-Free Efficiency for VLA Inference}

Training-free VLA efficiency methods reduce cost along three parts of the inference loop: the visual information given to the policy, the network computation used in each policy call, and how often the policy is invoked. We review prior work along these axes and position the five interventions evaluated in this study.

\noindent\textbf{Visual Efficiency.}
Classical foveated vision preserves high resolution near a region of interest and compresses the periphery~\cite{schwartz,traver}, and recent robot-learning methods use gaze-centered tokenization or gaze supervision to guide visual processing~\cite{lookfocusact,gazereg}. For a transformer VLA with a fixed token grid, however, changing the image in pixel space does not reduce computation, which motivates methods that select, prune, cache, or reuse visual representations. VLA-Cache reuses cached representations for stable patches and recomputes changed or task-relevant regions~\cite{vlacache}. On task-finetuned OpenVLA LIBERO checkpoints it reports $74.7\%$ average success against $75.0\%$ under dense inference with a $39\%$ CUDA-latency reduction, and it has become a reference for later training-free VLA acceleration methods~\cite{vlapruner,efficientvla,specprune}. VLA-Cache and VLA-Pruner also show that token-reduction methods transferred from generic VLMs, such as FastV and SparseVLM, can hurt VLA control because token importance for action prediction does not match generic visual-language saliency~\cite{fastv,sparsevlm,vlacache,vlapruner}.

\noindent\textbf{Network Depth Reduction.}
ShortGPT ranks layers by Block Influence, and Gromov et al.\ remove contiguous layer groups by representation similarity~\cite{shortgpt,gromov}. EfficientVLA applies non-contiguous layer pruning to VLAs without retraining~\cite{efficientvla}, MoLe-VLA learns an input-dependent layer router through self-distillation~\cite{molevla}, and compact VLA designs reduce depth at construction by using fewer language-model layers or by removing the language-model pathway from action prediction~\cite{flower,smolvla,turbovla}. The usefulness of depth reduction therefore depends on how much redundant depth a given backbone retains.

\noindent\textbf{Temporal Efficiency.}
Temporal redundancy is exploited by frame skipping and action chunking~\cite{dqn,act,diffusionpolicy,openvlaoft}. Skipping policy calls reduces inference cost but also changes how often the policy receives feedback, so recent methods use state-dependent signals. FlashVLA gates action reuse on visual and action consistency, and SpecPrune-VLA adapts computation to action dynamics~\cite{flashvla,specprune}. Other methods reuse temporal visual information without skipping policy calls. TTF-VLA fuses current and historical visual tokens by motion and attention relevance, VLA-InfoEntropy uses entropy-based signals, and VLA-IAP uses interaction alignment for visual-token selection~\cite{ttfvla,vlainfoentropy,vlaiap}.

\noindent\textbf{Static and Conditional Interventions.}
The interventions we evaluate divide by whether they use a signal. Visual foveation and fixed action repetition apply the same reduction throughout execution and serve as static controls. Depth pruning reduces computation according to a signal estimated once from calibration data, following the Block Influence criterion of ShortGPT~\cite{shortgpt}. Guarded reuse follows FlashVLA and SpecPrune-VLA in checking the newest observation, recent action consistency, gripper state, and translation before skipping a policy call~\cite{flashvla,specprune}. Temporal fusion reuses visual features across frames while still querying the policy at every step, so temporal visual redundancy can be studied separately from feedback skipping~\cite{ttfvla,vlainfoentropy,vlaiap}. Section~III specifies each mechanism.

\subsection{Evaluation of Efficient VLA Inference}

Efficient VLA methods are commonly evaluated on LIBERO and SimplerEnv~\cite{libero,simplerenv}, and recent efforts such as \texttt{vla-eval} and StarVLA aim to make evaluation consistent across models and codebases~\cite{vlaeval,starvla}. Acceleration studies still differ in their backbones, benchmarks, dense baselines, and inference backends, so it is difficult to tell whether a reported gain comes from the intervention itself or from the configuration in which it was tested. Multi-backbone studies often evaluate all models on one benchmark, while other works evaluate different models on different environment subsets~\cite{molevla,vlacache,specprune,vlapruner,gazereg,vlaiap}. Most results are aggregate success rates, which do not show whether the same episodes improve or fail after an intervention, and measured speedup depends on the dense implementation and attention backend used for reference~\cite{specprune}.

Our study addresses this gap by placing five representative training-free interventions, visual foveation, fixed action repetition, depth pruning, guarded action reuse, and temporal feature fusion, under one matched evaluation protocol. Each intervention is compared against the corresponding dense policy using the same task instances, matched episodes, and consistent runtime measurement. We exclude quantization because it changes numerical precision rather than visual, temporal, or network computation, and we exclude learned early-exit methods such as DeeR-VLA because they require additional training~\cite{deervla}. Following controlled ``bag-of-tricks'' evaluations~\cite{bagoftricks_cnn,bagoftricks_llm}, our goal is to compare representative training-free strategies under the same conditions rather than to introduce another isolated acceleration mechanism.
